% Generated from validated Article Draft Result. Do not hand-edit; revise the structured draft instead.
\begingroup\raggedright
\subsection{Agentic workloadはservingの前提を変える}
\par\endgroup

LLM servingを「独立したrequestが来てtokenを返す」問題としてだけ考えると、4月の研究群が狙った課題を捉えにくい。agentic workflowでは複数のLLM callが依存関係を持ち、同じcontextやKV stateを再利用し、toolの結果を待ちながら次のcallへ進む。multimodal MoEではroutingやload balanceの条件も変わる。4月は、こうした新しいworkloadを前提にscheduler、cache、power、workflow semanticsを見直すsystems researchと、実装stack側のvLLM更新が重なった。\autocite{src-0fc9d11483aa,src-2a6eae37bd39,src-551442a9e215,src-b790ec92aa01,src-d91793d0a6fa,src-e1e959857ced,src-fe12f92c05d7}


\begingroup\raggedright
\subsection{vLLM — model supportとschedulerの両方が動く}
\par\endgroup

vLLM v0.19.0は4月3日にreleaseされ、Gemma 4 architecture support、zero-bubble asynchronous scheduling with speculative decoding、multimodal／cache-management関連の変更をproject release notesに記録した。4月27日のv0.20.0ではinitial DeepSeek V4 supportと関連するserving/runtime変更が記録される。ここで重要なのは、model releaseへ追随するcompatibilityと、scheduler／cacheの内部実装が同じruntimeの更新として進む点である。どちらのrelease noteもprojectによる記録であり、記載されたsupportが全hardware・全workloadで同じ性能を保証するわけではない。\autocite{src-b790ec92aa01,src-fe12f92c05d7}


\begingroup\raggedright
\subsection{Scepsy — aggregate LLM pipelineをschedulerの単位にする}
\par\endgroup

Scepsyは4月16日の論文で、arbitrary multi-LLM agentic workflowをGPU clusterへ配置するserving systemとして提示された。従来の単一model requestではなく、複数LLMから構成されるworkflow全体をaggregate pipelineとして扱う発想が中心にある。これはagentの「賢さ」を改善する研究ではなく、agentが発生させる複数stageのinferenceをどのようにcluster上でscheduleするかというsystem問題である。著者が報告する効率改善は独立再現していないため、本稿では設計上の問題設定を主に扱う。\autocite{src-e1e959857ced}


\begingroup\raggedright
\subsection{KAIROSとPythia — contextとworkflow semanticsをservingへ持ち込む}
\par\endgroup

KAIROSは4月17日にcontext-aware power optimization system for agentic AI servingとして提案され、agentic workloadのcontextをpower管理へ利用する方向を示す。Pythiaは4月28日の初回投稿で、multi-agent workflowのsemanticsをserving layerへ渡すsimple interfaceを通じてoptimization機会を増やすと説明する。両者に共通するのは、serving layerをapplication semanticsから完全に切り離さず、agent側の状態やworkflow構造をoptimization signalとして利用しようとする点だ。ただしPythiaの保持locatorは5月14日のlater revisionであり、4月初稿にない記述を4月の成果として用いない。\autocite{src-551442a9e215,src-d91793d0a6fa}


\begingroup\raggedright
\subsection{CacheFlow — KV restorationを並列実行問題として扱う}
\par\endgroup

CacheFlowは4月28日に、KV cache restorationをmulti-dimensional parallel execution problemとして再定式化するframeworkを提案した。agentic／long-context workloadでは、cacheを持つか捨てるかだけでなく、保存済みstateをどこからどの並列度で復元するかがlatencyやresource利用に関わる。ここでの価値は個別benchmark数値より、KV cacheが単なるmemory optimizationからserving topologyの設計変数へ広がっているという問題設定にある。\autocite{src-0fc9d11483aa}


\begingroup\raggedright
\subsection{ReaLB — multimodal MoEのload balance}
\par\endgroup

ReaLBは4月21日の初回投稿で、multimodal MoE inference向けのreal-time load balancing methodを提案し、著者はzero scheduling overheadを主張する。multimodal inputではmodalityやtoken構成によってexpert routingの負荷が変わり得るため、MoE servingは単純な均等batch処理では捉えにくい。なお保持locatorは5月9日のlater revisionであるため、本号では初回投稿日をApril chronologyとして扱い、後日改訂された性能説明を4月へ遡及させない。\autocite{src-2a6eae37bd39}


